\subsection{The specified cubic obstruction with twisted coefficients}
\label{sec:cubic-twisted-obstruction}

The following construction identifies the actual cubic residue of
Theorem~\ref{thm:odd-torsion-cubic}. It supplies no arithmetic
nonvanishing assertion. In particular, the defining system below is fixed
by the actual class-two companions; it is not an arbitrary choice in an
ordinary triple Massey product.

\begin{proposition}[K17-CUBIC-TRIANGULAR]
\label{prop:cubic-triangular}
Use the actual weighted logarithmic modules of
Theorem~\ref{thm:odd-torsion-cubic}, with degrees \(1,2,3\) on
\(Q,T^+,T^-\), respectively. Put
\(R=\mathbb Z_p[[\Gamma]][1/p]\), write \(D\) for \(\gamma\), and suppose
\(T_a[1/p]\) is cyclic for a fixed odd character \(a\), with companion
polynomial \(q_a\). Fix a root \(\alpha\) of \(q_a\), a finite splitting
field \(L/\mathbb Q_p\) with integers \(\mathcal O\), and an ordered
triple of odd characters \(a_1+a_2+a_3=a\). Set
\[
 A=\mathcal O[[T_1,T_2]],\qquad
 z_1=1+T_1,\quad z_2=1+T_2,\quad
 z_3=\frac{\alpha}{z_1z_2},\qquad K=\operatorname{Frac}(A[1/p]).
 \tag{CT.1}
\]
There is a specified \(\Gamma\)-equivariant pushout of the actual
degree-three extension to a one-dimensional \(K\)-module on which
\(D\) acts as \(\alpha\). Its extension class is the actual cubic
companion numerator restricted to \(z_1z_2z_3=\alpha\). Thus its
nonsplitting is exactly the Triple-NV test for this ordered triple.
The assertion includes repeated roots and needs no cyclicity of the
other even or odd torsion components.
\end{proposition}

\begin{proof}
\par\medskip\noindent\textbf{The actual degree-two and degree-three modules.}
The degree-two module is the extension
\[
 0\longrightarrow N=\widehat{\bigwedge}^{2}Q[1/p]
 \longrightarrow E_2\longrightarrow T^+[1/p]\longrightarrow0.
 \tag{CT.2}
\]
The degree-three module is
\[
 0\longrightarrow W_3\longrightarrow E_3
 \longrightarrow T^-[1/p]\longrightarrow0,
 \qquad
 W_3=L_3(Q[1/p])\oplus[Q[1/p],T^+[1/p]]
 \tag{CT.3}
\]
as completed vector spaces. The displayed direct sum is not asserted to
split the \(\Gamma\)-action. Intrinsically, \(W_3\) is generated by
\([Q,E_2]\), with the Jacobi relations on three inputs from \(Q\).
The action on \(E_2\) determines its mixing of \([Q,T^+]\) into
\(L_3(Q)\).

Choose the rank-one \(R\)-basis in each odd character of \(Q[1/p]\).
Let \(f_i:Q_L\to K\) be evaluation at \(D=z_i\) on character
\(a_i\), and zero on the other characters. Then
\(f_i(Dx)=z_i f_i(x)\). The substitutions are integral and continuous
for maximal-ideal topologies: every root of the companion polynomials
lies in \(1+\mathfrak m_{\mathcal O}\). These coefficient rings have
bounded denominators; no identification with unrestricted
\(\mathbb Q_p[[T]]\) is used.

\par\medskip\noindent\textbf{The uniquely determined proper-pair coordinates.}
Define
\[
 n_{12}(x\wedge y)=f_1(x)f_2(y)-f_1(y)f_2(x),\qquad
 n_{23}(x\wedge y)=f_2(x)f_3(y)-f_2(y)f_3(x).
\]
Choose any nonzero integral polynomial \(q_+\) annihilating
\(T^+[1/p]\). Its characteristic polynomial on the finite integral
torsion lattice is one choice. Both
\(q_+(z_1z_2)\) and \(q_+(z_2z_3)=q_+(\alpha/z_1)\) are nonzero:
their arguments are nonconstant functions of the independent formal
variables. They are therefore invertible in \(K\). For \(y\in E_2\), put
\[
 b_{12}(y)=\frac{n_{12}(q_+(D)y)}{q_+(z_1z_2)},\qquad
 b_{23}(y)=\frac{n_{23}(q_+(D)y)}{q_+(z_2z_3)}.
 \tag{CT.4}
\]
The numerators are defined because \(q_+(D)y\in N\). These maps extend
the corresponding \(n_{ij}\) and have eigenvalues \(z_iz_j\).
An equivariant difference between two extensions factors through
\(T^+\); applying its annihilator, whose value is invertible, makes
that difference zero. This proves uniqueness and independence of the
chosen annihilator. On a cyclic even component, (CT.4) is precisely its
previously solved pair kernel. Pair-NV is not needed for this construction.

\par\medskip\noindent\textbf{The central functional and its extension class.}
In the strictly upper triangular Lie algebra of size four, send
\[
 x\longmapsto f_1(x)e_{12}+f_2(x)e_{23}+f_3(x)e_{34},\qquad
 y\longmapsto b_{12}(y)e_{13}+b_{23}(y)e_{24}.
\]
Their brackets on \([Q,Q]\) agree by the definition of \(n_{ij}\).
Consequently the \((1,4)\)-entry of the matrix bracket defines
\[
 \psi([x,y])=f_1(x)b_{23}(y)-b_{12}(y)f_3(x),\qquad
 \psi:W_3\longrightarrow K.
 \tag{CT.5}
\]
Jacobi holds because it holds for matrix commutators. This retains both
mixed contributions from the actual even extension. The map has
\(\Gamma\)-eigenvalue \(z_1z_2z_3=\alpha\) and \(\Delta\)-character
\(a\). Push out (CT.3) in character \(a\) along \(\psi\):
\[
 0\longrightarrow K_\alpha\longrightarrow E_{a,\psi}
 \longrightarrow T_a[1/p]\longrightarrow0.
 \tag{CT.6}
\]
Here \(R_L=\mathcal O[[\Gamma]][1/p]\) acts on \(K_\alpha\) through
\(D\mapsto\alpha\). The two-term free resolution of
\(R_L/(q_a)\) gives
\[
 \operatorname{Ext}^1_{R_L}(R_L/(q_a),K_\alpha)
   =K/q_a(\alpha)K=K.
 \tag{CT.7}
\]
If \(v\) is the chosen cyclic vector and \(y\) is its actual
degree-three lift, the class in (CT.7) is
\(\psi(q_a(D)y)\): lift the free generator to \(y\) and apply the
relation \(q_a\). A change of lift adds \(q_a(\alpha)\) times a
scalar, hence adds zero. This computation uses the full companion
polynomial, including its multiplicities; no derivative is taken.
Substitution of (CT.4) into (CT.5) is exactly the substitution of the
solved pair kernels into the actual odd companion error. Thus
\(\psi(q_a(D)y)\) is the stated cubic residue.

There is also a Lie lifting interpretation. Modulo its central
\((1,4)\)-entry, the above map is a fixed equivariant map to the
two-step quotient of the triangular Lie algebra. A lift over the odd
torsion component exists exactly when (CT.6) splits. The remaining
odd quotient \(Q_L\) is free of finite rank over \(R_L\), so its
module extension by the central degree-three part splits by choosing
lifts of finitely many basis elements. Choose these in the rational
compact Lie module and use its continuous \(\Lambda\)-action; the
section has one bounded \(p\)-denominator. Changing those lifts by
degree-three elements changes no bracket through degree three.
It follows that this free quotient contributes no additional
obstruction to the specified triangular Lie lift.
\end{proof}

\begin{proposition}[K17-CUBIC-TWISTED]
\label{prop:cubic-twisted}
In the notation of Proposition~\ref{prop:cubic-triangular}, put
\(P=H\rtimes(\Gamma\times\Delta)\). After explicit nonzero scalar
rescalings which preserve the test over \(K\), the specified partial
system is defined over the compact ring \(A\). It has continuous
first-coordinate cocycles \(x_i\) with coefficient characters
\[
 \chi_i(\gamma)=z_i,\qquad
 \chi_i(\delta)=\omega(\delta)^{a_i},\qquad
 \chi_i|_H=1.
 \tag{CT.8}
\]
There are specified continuous pair cochains \(B_{12},B_{23}\) such that
\[
 dB_{12}=-x_1\smile x_2,\qquad dB_{23}=-x_2\smile x_3.
\]
The continuous cocycle
\[
 c=x_1\smile B_{23}+B_{12}\smile x_3
 \quad\hbox{with coefficients in }A_{\chi_1\chi_2\chi_3}
 \tag{CT.9}
\]
is the obstruction to the specified twisted triangular lift. More
precisely,
\[
 [c]\otimes1\ne0
 \quad\hbox{in}\quad
 H^2\bigl(C^\bullet_{\rm cts}(P,A_{\chi_1\chi_2\chi_3})\bigr)
       \otimes_A K
 \quad\Longleftrightarrow\quad
 \psi(q_a(D)y)\ne0.
 \tag{CT.10}
\]
Only algebraic localization of the displayed continuous cochain complex
is used. It is not identified with the complex of all continuous
\(K\)-valued cochains.
\end{proposition}

\begin{proof}
\par\medskip\noindent\textbf{Compact coefficients before cohomology.}
Let \(d_{12}=q_+(z_1z_2)\) and \(d_{23}=q_+(z_2z_3)\), nonzero
elements of \(A\). Replace \(f_1,f_2,f_3\) by
\(d_{12}f_1,f_2,d_{23}f_3\). By (CT.4), the new pair maps are
\(d_{12}b_{12}\) and \(d_{23}b_{23}\); their proper-pair
denominators cancel. The new central functional is
\(d_{12}d_{23}\psi\). All these factors are units in \(K\).
The rational coordinate maps on the finitely generated
\(\Lambda\)-module \(Q\) have bounded \(p\)-denominators.
Their completed bilinear maps are bounded, and the finite torsion
bases add only finitely many bounds. Multiply every first coordinate
by a common \(p^N\). The pair and central coordinates are multiplied
by \(p^{2N}\) and \(p^{3N}\), respectively. Increasing \(N\)
absorbs the finitely many degree-three logarithm and exponential
denominators, also when \(p=3\). Thus the partial system and central
relation can be represented over \(A\). In what follows the same
letters denote these rescaled maps; their zero tests over \(K\)
are unchanged.

\par\medskip\noindent\textbf{The group coordinates and their twists.}
Exponentiate the specified Lie map modulo its central entry. For
\(h\in H\), let \(x_i(h)\) be its first coordinates and let
\(B_{12}(h),B_{23}(h)\) be its group coordinates in positions
\((1,3),(2,4)\). In particular, the first pair coordinate is the
Lie coordinate plus \(x_1(h)x_2(h)/2\); the second has the
analogous correction. Define them on \(h t\),
\(t\in\Gamma\times\Delta\), by their values on \(h\).
The diagonal matrix
\[
 D_0(t)=\operatorname{diag}
   (\chi_1\chi_2\chi_3(t),\chi_2\chi_3(t),\chi_3(t),1)
\]
implements the action on the triangular entries. The resulting
map to the upper triangular group modulo its normal
\((1,4)\)-subgroup is a homomorphism, equivalently a twisted
unipotent cocycle. This subgroup is central in the unipotent
group; the diagonal group acts on it by
\(\chi_1\chi_2\chi_3\).

For the inhomogeneous differential
\(dF(g,h)=gF(h)-F(gh)+F(g)\), its first coordinates satisfy
\(dx_i=0\). Matrix multiplication gives
\[
 \begin{split}
 B_{12}(gh)&=B_{12}(g)+\chi_1\chi_2(g)B_{12}(h)
                      +x_1(g)\chi_2(g)x_2(h),\\
 B_{23}(gh)&=B_{23}(g)+\chi_2\chi_3(g)B_{23}(h)
                      +x_2(g)\chi_3(g)x_3(h).
 \end{split}
 \tag{CT.11}
\]
These are the asserted pair differentials. The cup convention is
\((F\smile G)(g,h)=F(g)gG(h)\). Its Leibniz rule shows directly
that \(dc=0\): the two triple products cancel. In entries,
\[
 c(g,h)=x_1(g)\chi_2\chi_3(g)B_{23}(h)
              +B_{12}(g)\chi_3(g)x_3(h).
 \tag{CT.12}
\]
A central coordinate \(Z(g)\) gives a full lift exactly when
\(dZ=-c\). This proves the obstruction formula, including its
coefficient action.

There is an ordinary connecting homomorphism with this exact value.
Put
\[
 L(g)=\chi_2\chi_3(g)
       \begin{pmatrix}\chi_1(g)&x_1(g)\\0&1\end{pmatrix},
 \qquad
 R_0(g)=\begin{pmatrix}\chi_3(g)&x_3(g)\\0&1\end{pmatrix}.
\]
Both are matrix representations. Give \(V=M_2(A)\) the action
\(gX=L(g)X R_0(g)^{-1}\). Its top-right line is
\(A_{\chi_1\chi_2\chi_3}\). Write \(V'\) for the quotient
by that line. The actual cross-block of the partial matrix is
\[
 U(g)=\begin{pmatrix}
             B_{12}(g)\chi_3(g)&0\\
             x_2(g)\chi_3(g)&B_{23}(g)
       \end{pmatrix}.
\]
The class of \(U(g)R_0(g)^{-1}\) in \(V'\) is a
continuous cocycle \(\kappa'\). Its indicated lift
\(\widetilde\kappa=UR_0^{-1}\) satisfies
\[
 d\widetilde\kappa(g,h)
  =\bigl(L(g)U(h)+U(g)R_0(h)-U(gh)\bigr)R_0(gh)^{-1}
  =c(g,h)e_{12}.
 \tag{CT.13}
\]
The last equality uses (CT.11); the right multiplication fixes the
top-right matrix unit. Hence the connecting map for
\(0\to A_{\chi_1\chi_2\chi_3}\to V\to V'\to0\)
sends \([\kappa']\) to \([c]\). This is a specified finite
coefficient extension, not an assumed identification with a full
augmentation-graded Iwasawa module.

\par\medskip\noindent\textbf{The bounded lift equivalence.}
The ring \(K\) is the algebraic localization of \(A\) at its
nonzero elements. Flatness gives
\[
 H^2\bigl(C^\bullet_{\rm cts}(P,A_\chi)\otimes_A K\bigr)
   =H^2\bigl(C^\bullet_{\rm cts}(P,A_\chi)\bigr)\otimes_A K,
 \qquad \chi=\chi_1\chi_2\chi_3.
 \tag{CT.14}
\]
This assertion is solely an identity for localization of a complex
of \(A\)-modules. A cochain on the left has one common scalar
denominator; no topology on all \(K\)-valued cochains is invoked.

Suppose the left side of (CT.10) is zero. There is a nonzero
\(s\in A\) and a continuous \(A\)-valued cochain \(v\) with
\(dv=sc\). Then \(Z=-v/s\) supplies the missing central
coordinate. It is a bounded scalar fraction of a continuous cochain.
The resulting full triangular homomorphism on \(P\) is
\(\Gamma\times\Delta\)-equivariant on \(H\): on the outer
subgroup its unipotent part is central, and so its conjugation on
\(H\) is the prescribed diagonal action.

Take logarithms through degree three. This operation uses finitely
many rational coefficients. It respects the assigned filtration:
the fixed first coordinates kill all torsion; an even torsion
eigenlift has no first coordinate and its logarithm has no central
odd-character coordinate; an odd torsion eigenlift has no first
or second coordinate. Thus their respective starting degrees are
at least \(1,2,3\). Substitution into the weighted Magnus
expansion consequently gives a Lie map on the actual quotient
through assigned degree three. This also holds for the infinite
compact basis: its lifts tend to one, the coefficient maps are
continuous with a common denominator, and only the first three
completed tensor degrees occur. Clearing that denominator and
the finite logarithm denominators gives continuous maps on those
compact tensor modules. The restriction to \(E_3^{[a]}\) is an
equivariant extension of \(\psi\), so (CT.6) splits. By (CT.7)
the cubic residue is zero.

Conversely, suppose \(\psi(q_a(D)y)=0\). Choose the actual
conjugate lifts \(y,Dy,\ldots,D^{r-1}y\), where \(r=\deg q_a\).
Their classes form a rational basis of \(T_a\). Set their
central values to zero and use \(\psi\) on \(W_3\). The final
companion relation is compatible precisely because its value is
\(\psi(q_a(D)y)=0\). This gives the equivariant extension on
\(E_3^{[a]}\); it is zero on the other character components.
Its coefficients have a common denominator: only a finite basis
change on the integral torsion lattice is involved. For the
remaining odd quotient \(Q_L\), choose the finitely many
\(R_L\)-basis lifts as in the preceding proposition. Their
continuous \(\Lambda\)-spans have one bounded \(p\)-denominator.
Consequently the full Lie lift, and its exponential through degree
three, have one common denominator on the actual compact source.
Its central coordinate \(Z\) therefore belongs to
\(C^1_{\rm cts}(P,A_\chi)\otimes_A K\) and satisfies \(dZ=-c\).
This proves the reverse implication and (CT.10). The construction
retains the full cyclic companion, so the same proof covers repeated
roots without a simple-root specialization.
\end{proof}

\begin{proposition}[K17-CUBIC-CYCLOTOMIC-INPUTS]
\label{prop:cubic-cyclotomic-inputs}
In the arithmetic setting, replacing the three abstract rank-one
linear evaluations by those supplied by the actual norm-compatible
cyclotomic \(p\)-unit family changes the test in (CT.10) by a
nonzero element of \(K\).
\end{proposition}

\begin{proof}
For odd \(a_i\), put \(\theta_i=\omega^{1-a_i}\), an even
character. At \(F_0=\mathbb Q(\mu_p)\), the \(\theta_i\)-projection
of \(1-\zeta_p\) is nonzero in the rational \(p\)-unit space.
For nontrivial \(\theta_i\), the cyclotomic units span its real
unit line by Hasse's finite-index theorem. For the trivial character,
the norm is \(p\), so its valuation direction is nonzero. The
precise primary locators are \cite[Theorem~4.3.3 and its proof,
pp.~108--109; Proposition~6.1.5 and its proof,
p.~152]{SharifiIwasawaFull2026}. No integral unit-index assertion or
Vandiver assumption is used here.

The Kummer class is therefore non-torsion. The restriction kernel
from \(H^1(P_0,\mathbb Z_p(1))\) to \(H^1(H,\mathbb Z_p(1))\)
is inflated from \(H^1(\Gamma,\mathbb Z_p(1))\). Since
\(\Delta\) commutes with \(\Gamma\), this inflated group has
only the odd cyclotomic \(\Delta\)-character. It cannot kill the
even Kummer component. The explicit norm-transfer map of
K17-NORM-TORSION is consequently nonzero on \(Q_{a_i}\); by that
lemma it kills all \(\Lambda\)-torsion.

After the involution \(\gamma\mapsto u\gamma^{-1}\), this
nonzero rank-one map is a nonzero bounded one-variable multiplier
\(h_i\) times \(f_i\). The notation \(h_i(z_i)\) means evaluation
of that power series in \(z_i-1\), after the indicated involution.
The pair maps acquire factors \(h_1(z_1)h_2(z_2)\) and
\(h_2(z_2)h_3(z_3)\), and the central class acquires
\(h_1(z_1)h_2(z_2)h_3(z_3)\). Each is nonzero in \(K\).
For the third factor, substitution is injective by the integral
coordinate change \(V=z_3-1,W=T_1\), whose inverse is
\(z_2=\alpha/(z_1(1+V))\); all constant shifts lie in the
maximal ideal. Hence the total factor is a unit in \(K\), proving
the assertion.

The global cyclotomic cocycles can also be represented by the
zero-on-\(\Gamma\) first-coordinate cocycles used above after the
same algebraic localization. Indeed, after their restrictions to
\(H\) have been matched by these multipliers, their difference
\(\eta\) is inflated from \(\Gamma\times\Delta\). Since
\(z_i-1\ne0\) in \(A\), subtract the coboundary of
\(\eta(\gamma)/(z_i-1)\). The resulting cocycle vanishes on
\(\Gamma\) by continuity. For \(\delta\in\Delta\), the
cocycle identity for the commuting elements \(\gamma,\delta\)
then gives \((z_i-1)\eta(\delta)=0\), so it vanishes on
\(\Delta\) as well. This uses a single additional scalar
denominator and identifies the localized \(H^1\) input classes.
It does not replace the specified pair defining system by arbitrary
pair primitives: those remain the actual coordinates (CT.4) and
(CT.11).
\end{proof}

\paragraph{Exact source comparison and the remaining arithmetic gate.}
The general twisted matrix definition is already covered by
\cite[Definitions~3.1.3, 3.2.1--3.2.2, pp.~22--23]{LLSWW2023}.
Their subsequent proper-system and Bockstein comparison specializes
to a single possibly nontrivial coefficient module, with untwisted
outer first coordinates
\cite[Definitions~3.3.1--3.3.2 and Theorems~3.3.4, 4.1.2,
pp.~24--29]{LLSWW2023}. It cannot be applied verbatim to the
three distinct universal characters in (CT.8). Equations
(CT.11)--(CT.14) prove the needed twisted cochain identity directly;
they do not identify its finite coefficient extension with an
augmentation sequence whose full cokernel is known to vanish.

In the universal cohomological variables \(w_i=u/z_i\), the
triple divisor is \(w_1w_2w_3=u^3/\alpha\), and its dual
character is \(\omega^{3-a}\). To obtain actual Triple-NV one
must still prove that the specified class (CT.9) has nonzero image
in (CT.10) for every required odd companion root, using some
matching ordered triple for each root. A full Bockstein-cokernel
identity without the corresponding image-support or cokernel
vanishing statement does not prove this. Nor does coherent
vanishing in a repeated cyclotomic direction decide this different
three-input defining system. The propositions above establish the
actual class and its bounded coefficient realization; they assert
no new arithmetic Triple-NV case and no new full-kernel theorem.
